\section{What is known, measured, and unknown}
%==========================================================================

\begin{adjustwidth}{-1.4cm}{-1.4cm}
\begin{table}[htbp]
\centering
\footnotesize
\caption{Epistemic status of each quantity. The \emph{Role} column gives the
system-theoretic classification; the \emph{Status} column gives what is actually
available from clinical records.}
\label{tab:status}
\begin{tabular}{@{}p{0.085\textwidth} p{0.175\textwidth} p{0.135\textwidth} p{0.20\textwidth} p{0.30\textwidth}@{}}
\toprule
\textbf{Quantity} & \textbf{Meaning} & \textbf{Role} & \textbf{Status} & \textbf{Source and caveat} \\
\midrule
$m(t)$ & plasma 3-hydroxypropionate, the marker the state denotes &
state; the observable the model is written in &
measured, $\approx 6$--$7$ samples/year &
a $50\%$ lower 3-HP is associated with a decompensation rate ratio of $0.7$
($p<0.01$)~\cite{natural_history}. % C19
\emph{Proxy}: the enzyme substrate is intracellular propionyl-CoA, never measured \\
\addlinespace
C3 & propionylcarnitine &
a second marker, \textbf{not} interchangeable with $m$ &
measured alongside 3-HP &
a $50\%$ lower C3 gives a rate ratio of $0.6$ ($p=0.017$)~\cite{natural_history}: % C19
carries risk information, but on its own scale, so it is not a substitute for 3-HP \\
\addlinespace
2-MC & 2-methylcitrate &
measured, but \textbf{not} used as the state &
measured alongside 3-HP &
rate ratio $1.0$ ($p=0.84$)~\cite{natural_history}: % C19
no detectable association with decompensation. Excluded from the state on that ground \\
\addlinespace
$s(t)$ & propiogenic substrate load &
exogenous disturbance, partly structured &
partially known &
prescribed protein intake and illness dates recorded; episode amplitudes unknown;
unrecorded triggers exist \\
\addlinespace
$\theta$ & residual PCC capacity: the Michaelis-Menten $V_{\max}$ &
\textbf{parameter, patient-specific, constant} &
\textbf{unknown} &
set by genotype; measurable \emph{ex vivo} in fibroblasts or lymphocytes, but the
in vitro to in vivo mapping is uncertain \\
\addlinespace
$K_m$ & Michaelis constant of PCC: the level at which the enzyme runs at half of
$V_{\max}$ &
parameter, assumed population-level &
unknown &
property of the active site, assumed genotype-independent \\
\addlinespace
$k_e$ & non-enzymatic clearance; a rate constant in Model A, the $V_{\max}$ of the
alternative route in Model B &
parameter, assumed population-level &
unknown &
carnitine conjugation and renal excretion. \emph{Caveat}: carnitine is also administered
therapeutically, so partly clinician-manipulated \\
\addlinespace
$u_c(t)$ & carnitine dose &
\textbf{input on $k_e$}, selective but \emph{endogenous}; the working perturbation of
this note &
recorded in charts, present in every patient &
adjusted in response to patient state, hence confounded by indication. Quasi-exogenous
instruments: weight-based recalculation after growth, supply interruptions, changes of
formulation \\
\addlinespace
$\Delta(t)$ & therapeutic increment to capacity &
\textbf{input on $\theta$}, selective and exogenous &
\textbf{not administered in the accessible cohort} &
investigational. Structurally interchangeable with $u_c$: by
Proposition~\ref{prop:channel} the determinant of the identification problem is the same
for both channels, so the analysis is carried out on $u_c$ and transfers to $\Delta$.
Decay rate $\lambda$ possibly patient-specific \\
\addlinespace
$\bar m$ & tolerable metabolite level &
constant in the definition of $\thc$ &
\textbf{not measured} &
clinical convention, not a physiological constant \\
\addlinespace
$\spk$ & peak substrate load in illness &
extreme value of the disturbance &
\textbf{not measured} &
varies with severity of the intercurrent event \\
\addlinespace
$N(t)$ & decompensation events &
counted output, threshold-triggered &
counted and dated &
admissions reliably coded; only a minority of all events lead to admission \\
\bottomrule
\end{tabular}
\end{table}
\end{adjustwidth}

\begin{adjustwidth}{-1.4cm}{-1.4cm}
\begin{table}[htbp]
\centering
\footnotesize
\caption{The three candidate models. They are alternative hypotheses about the
non-enzymatic clearance term, not stages of a development.}
\label{tab:models}
\begin{tabular}{@{}p{0.035\textwidth} p{0.235\textwidth} p{0.175\textwidth} p{0.185\textwidth} p{0.145\textwidth} p{0.13\textwidth}@{}}
\toprule
& \textbf{Clearance structure} & \textbf{Equilibria} & \textbf{Question answered}
& \textbf{Cost} & \textbf{Status} \\
\midrule
\textbf{A} &
$\theta g(m) + k_e m$ \newline Michaelis-Menten enzyme, second term linear and unbounded &
one, stable, never lost &
Can $\theta$ be identified from clinical data? &
predicts no crisis at all &
working model for identifiability \\
\addlinespace
\textbf{B} &
$\theta g(m) + k_e h(m)$ \newline both Michaelis-Menten, ceiling $\theta + k_e$ &
one stable, lost when load exceeds the ceiling &
Under what condition is a tolerable steady state lost? &
both terms share a functional form, weakening separability &
required for $\thc$; not verified \\
\addlinespace
\textbf{C} &
$\theta g(m) + k\,c(t)\,m$ \newline with $c$ depleted by elimination &
two, stable and unstable, colliding at a saddle-node &
Why does recovery require active reduction of the load below baseline, rather than mere
removal of the trigger? &
adds an unmeasured state $c$ and an unknown $R$ &
hypothesis; discriminating test available in records. The only one of the three in which
$\thc$ is well posed \\
\bottomrule
\end{tabular}
\end{table}
\end{adjustwidth}

\begin{remark}[What each model is for]
Models A, B and C are not stages of a development, and they are not ranked by realism.
They are alternative hypotheses about one term, and each is present for a different
reason.

Model A is used for the identifiability analysis because it is the version in which the
two clearance channels have maximally different functional forms, and is therefore the
\emph{most favourable} case rather than the most accurate. Its known defect is that it
predicts no crisis.

Model B is what is required to say anything about the loss of a tolerable steady state,
and it pays for that in identifiability, since both channels acquire the same functional
form.

Model C is usually introduced as the most biologically faithful of the three, which
undersells it. Its distinguishing property is that it is \textbf{the only one of the three
in which $\thc$ is well posed}. In A and B the critical capacity must be defined by
convention through $\bar m$; in C it is derived, as the point at which the attractor and
the separatrix merge at a saddle-node. See Proposition~\ref{prop:intrinsic}.
\end{remark}

\vspace{0.5em}
\hrule
\vspace{1em}

%==========================================================================
